%! ~~~ Packages Setup ~~~ 
\documentclass[]{../../../../tex/classes/styledArticle}
\usepackage{../../../../tex/packages/hebrewSupport}
\usepackage{../../../../tex/packages/mathShortcuts}
\usepackage{../../../../tex/packages/theoremsSupport}
\usepackage{../../../../tex/packages/enfancysections}

%! ~~~ Document ~~~

\author{שחר פרץ}
\title{\textit{לינארית 2א 8}}
\date{5 במאי 2025}
\begin{document}
    \maketitle
    \section{\en{T-Invariant spaces}}
    בגדול סיימנו את הפרק על חוגים. להלן תוספת קטנה מקורס לינארית 2 של האונ' העברית, שנמצאת באופן חלקי בחומר שלנו. 
    \defi{נניח ש־$V$ מ''ו מעל $\F$, ו־$T \co V \to V$ ט''ל. אז $U \subseteq V$ תמ''ו נקרא $T$־אינווריאנטי/$T$-שמורה אם לכל $u \in U$ מתקיים $T(u) \in U$. }
    \textbf{דוגמאות. }$V, \{0\}$ הם $T$־אינווריאנטים. גם המ''ע (המרחבים העצמיים) הם $T$־אינוו'. 
    
    \noti{שימו לב: אם $U \subseteq V$ תמ''ו אינווריאנטי, אז $T|_U \co U \to U$ ט''ל. }
    
    נניח ש־$u_1 \dots u_k$ בסיס ל־$U$ כנ''ל, ו־$W \subseteq V$ תמ''ו כך ש־$U \oplus W = V$, ונגיד ש־$w_{k + 1} \dot w_n$ בסיס ל־$W$, אז $B = (u_1 \dots u_k, w_{k + 1} \dots w_n)$ מקיים: 
    \[ [T]_B = \pms{[T_{|_U}] & A \\ 0 & B} \]
    (כאשר $[T_{|_U}] \in M_{k}$ ו־$B \in M_k$). 
    
    \defi{יהי $V$ מ''ו מעל $\F$, \ $T \co V \to V$ ט''ל ו־$v \in V$ וקטור. אז \textit{תת־המרחב־הציקלי} הנוצר מ־$T$ ע''י $T$ הוא 
    \[ \mathcal{Z}(T, v) := \Sp\{T^k(v) \mid k \in \N_0\} \]}
    \textbf{טענה. }
    \begin{itemize}
        \item $\zc(T, v)$ תמ''ו של $V$ – טרוויאלי .
        \item $\zc(T, v)$ תמ''ו $T$־־אינ' – טרוויאלי גם. 
    \end{itemize}
    
    עתה נציג משהו נחמד. אם $V$ נוצר סופית, גם $\zc(T, v)$ נ''ס. נגיד שיהיה $n \in \N_0$ מינימלי, כך ש־$\zc(T, v) = \Sp\{v, Tv, \dots T^{n - 1} v\}$. אז $T^nv \in \zc(T, v)$. לכן קיימים $a_0 \dots a_{n - 1}$ כך ש־$T^nv + a_{n - 1}T^{n - 1}v + \cdots + a_0v = 0$. ניתן לקחת כבסיס את $v, Tv \dots T^{n - 1} v$ של $\zc(T, v)$. אז: 
    \[ [T]_B = \pms{0 & 0 & \cdots & 0 & -a_0
    \\ 1 & 0 & \cdots & 0&-a_1
    \\ 0 & 1 & \vdots & 0& \vdots
    \\ 0 & 0  & \cdots & 1 &-a_{n - 1}} \]
    כאשר השורה האחרונה כי: 
    \[ T(T^{n - 1}v)  = -\sum_{k = 0}^{n - 1}a_kT^kv \]
    כלומר, $A_f = [T]_B$ היא המטריצה המצורפת לפולינום $f(x) = x^{n} + a_{n - 1}x^{n - 1} + \cdots + a_0$. 
    
    \section{\en{Extending Fields}}
    \defi{$I \subseteq R$ אידיאל נקרא \textit{ראשוני} אם $\forall a, b \in R \co (a \cdot b) \subseteq P$, אז $(a) \in P \lor (b) \in P$. }
    \defi{אידיאל $I \subseteq R$ נקרא \textit{אי־פריק} אם $\forall a, b \in R$ אם $I = (a \cdot b)$ אז $I = (a)$ או $I = (b)$. }
    ראינו, שבתחום ראשי אי פריק=ראשוני. ניתן להראות באופן שקול כי: 
    \theo{$R$ תחום ראשי, אז $I$ ראשוני אמ''מ $I$ אי־פריק. }
    \defi{יהי $R$ תחום שלמות [אפשר להתעסק גם עם אידיאל ימני ושמאלי] ונניח ש־$I \subseteq R$ אידיאל. אז $R/_I := \{a + I \mid a \in R\}$ הוא חוג (בהגדרת $a + I = \{a + i \mid i \in I\}$ חיבור מנות), כאשר הפעולות: 
    \begin{itemize}
        \item $(a + I) + (b + I) = (a + b) + I$
        \item $(a + I)(b + I) = ab + I$
    \end{itemize}}
    צריך להוכיח שזה לא תלוי בנציגים (הנציגים $a, b$) והכל אבל בן לא עומד לעשות את זה. 
    \theo{בתחום ראשי $R$, אם $I$ אידיאל אי־פריק, אז $R/_I$ שדה. }
    
    \textbf{דוגמאות. }
    \begin{itemize}
        \item $\Z/_{\la p \ra}$ שדה. 
        \item $\R[x]$ תחום ראשי, ידוע $x^2 + 1$ אי־פריק. לכן $\R[x]/_{\la x^2 + 1\ra}\cong \C$. הרעיון: נוכל להסתכל על $p$ פולינום המבוטא כמו: 
        \[ p(x) = q(x) \cdot (x^2 + 1) + ax + b = ax +b + \la x^2 + 1\ra \]
        ואם נכפיל שני יצורים כאלו: 
        \[ ((ax + b) + I) + (cx + d + I)) = acx^2 + (ad + bc)x + bd + I \]
        אך ידוע ש־$x^2 + 1 = 0$ (כי זה האידיאל שלנו) עד לכדי נציג, כלומר מתקיים שוויון ל־$bd - ac + (ad + bc) + I$ זהו כפל מרוכבים. אפשר גם להראות קיום הופכי וכו'. 
    \end{itemize}
    \begin{proof}
        יהי $0 \neq a + I \in R/_I$. אם $a \neq 0$, אז ב־$R$ מתקיים $p \nmid a$ אי $p$ א''פ (אם הוא היה מחלק את $a$ אז $a = 0$) ולכן $p, a$ זרים (כי האידיאל אי פריק וכו'). אז קיימים $r, s \in R$ כך ש־$ar + ps = 1$. סה''כ: 
        \[ ar + ps + I = 1 + I \implies ar + I = 1 + I \quad \top \]
        לכן $r + I$ הופכי של $a + I$ וסיימנו. 
    \end{proof}
    (למעשה זה אממ). 
    
    
    \ndoc
\end{document}
